\documentclass{ximera}
\title{Functions}
\begin{document}
\begin{abstract}
What a function is (domain, codomain, range), the standard families of functions, even and odd functions, increasing and decreasing functions, and how to combine and compose functions. Textbook §1.1–1.3.
\end{abstract}
\maketitle
A function is nothing more than a rule where you give it something and get back something else. As an example:
\begin{image}
\includegraphics[width=0.4\textwidth]{b48a5324-93dc-45ee-b7c2-dc0b5706bda6-004_161_273_612_915}
\end{image}

More explicitly, a function $f$ is a rule that assigns to each $x$ in a set $D$ to exactly one element $\answer{f(x)}$, in a set $E$. This is normally denoted
\wordChoice{\choice[correct]{$f\colon D \rightarrow E$}\choice{$f\colon E \rightarrow D$}\choice{$D\colon f \rightarrow E$}}
where $D$ is called the domain and $E$ is called the codomain. The range of $f$ is the set of all possible values of $f(x)$ for every $x$ in $D$.

\begin{exercise}\label{exercise:1-1}
Create your own function that does not use any numbers nor atoms!
\begin{freeResponse}\end{freeResponse}
\end{exercise}

\begin{example}\label{example:1-2}
Let's look at some examples.
\begin{enumerate}
  \item A linear function is a function whose graph is a line (e.g. functions of the form $f(x)=m x+b$. Let $f$ be the linear function $f: \mathbb{R} \rightarrow \mathbb{R}$ such that
  % AUTHOR CHOICE: example chosen during conversion, change freely
  $f(x)=2x+1$.

  Domain: $\answer{\mathbb{R}}$ \qquad Codomain: $\answer{\mathbb{R}}$ \qquad Range: $\answer{\mathbb{R}}$
  \begin{image}
  \includegraphics[width=0.5\textwidth]{b48a5324-93dc-45ee-b7c2-dc0b5706bda6-004_361_378_1981_915}
  \end{image}

  \item A polynomial $f$ is a function of the form
  \[
  f(x)=a_{n} x^{n}+a_{n-1} x^{n-1}+\cdots+a_{2} x^{2}+a_{1} x^{1}+a_{0}
  \]
  where $a_{0}, \ldots, a_{n}$ are known as the coefficients and, if $a_{n} \neq 0$, then $f$ is said to have degree $n$.

  Let $f: \mathbb{R} \rightarrow \mathbb{R}$ be the polynomial of degree 2 such that
  % AUTHOR CHOICE: example chosen during conversion, change freely
  $f(x)=x^{2}$.

  Domain: $\answer{\mathbb{R}}$

  Codomain: $\answer{\mathbb{R}}$

  Range: $\answer{[0,\infty)}$
  \begin{image}
  \includegraphics[width=0.5\textwidth]{b48a5324-93dc-45ee-b7c2-dc0b5706bda6-005_363_378_834_958}
  \end{image}

  \item Let $f$ be the function $f: \mathbb{R} \rightarrow \mathbb{R}$ such that
  % AUTHOR CHOICE: example chosen during conversion, change freely
  $f(x)=x^{3}$.

  Domain: $\answer{\mathbb{R}}$

  Codomain: $\answer{\mathbb{R}}$

  Range: $\answer{\mathbb{R}}$
  \begin{image}
  \includegraphics[width=0.5\textwidth]{b48a5324-93dc-45ee-b7c2-dc0b5706bda6-005_363_378_1377_915}
  \end{image}

  \item A rational function $f$ is a ratio of two polynomials:
  \[
  f(x)=\frac{P(x)}{Q(x)} .
  \]
  Let $f$ be the rational function
  % AUTHOR CHOICE: example chosen during conversion, change freely
  \[
  f(x)=\frac{1}{x}.
  \]
  Domain: $\answer{(-\infty,0)\cup(0,\infty)}$ \qquad Range: $\answer{(-\infty,0)\cup(0,\infty)}$

  Why am I not asking for the codomain?
  % TODO: original had a blank here; intended answer unclear (the lecturer answers this live; probably
  % "because the domain is no longer all of R, and the codomain is just whatever set we choose, e.g. R")
  \begin{freeResponse}\end{freeResponse}
  \begin{image}
  \includegraphics[width=0.5\textwidth]{b48a5324-93dc-45ee-b7c2-dc0b5706bda6-006_361_378_799_958}
  \end{image}

  \item Piecewise defined function: is a function which has different formulas in different parts of their domains. Let $f$ be the piecewise defined function $f: \mathbb{R} \rightarrow \mathbb{R}$ such that
  % AUTHOR CHOICE: example chosen during conversion, change freely
  \[
  f(x)=\begin{cases} 1-x & \text{if } x \leq 1 \\ x^{2} & \text{if } x>1 \end{cases}
  \]
  % work: for x<=1, 1-x takes all values in [0,inf); for x>1, x^2 takes (1,inf); union is [0,inf)
  Domain: $\answer{\mathbb{R}}$ \qquad Codomain: $\answer{\mathbb{R}}$ \qquad Range: $\answer{[0,\infty)}$
  \begin{image}
  \includegraphics[width=0.5\textwidth]{b48a5324-93dc-45ee-b7c2-dc0b5706bda6-006_363_378_1571_915}
  \end{image}
\end{enumerate}
\end{example}

An \wordChoice{\choice[correct]{even}\choice{odd}\choice{increasing}} function is a function $f$ that satisfies $f(-x)=f(x)$. An \wordChoice{\choice[correct]{odd}\choice{even}\choice{decreasing}} function is a function $f$ that satisfies $f(-x)=-f(x)$.

\begin{example}\label{example:1-3}
Let $f: \mathbb{R} \rightarrow \mathbb{R}$ be the function
% AUTHOR CHOICE: example chosen during conversion, change freely
$f(x)=x^{3}$. Is it even or odd? \wordChoice{\choice{even}\choice[correct]{odd}\choice{neither}}
\end{example}

\begin{example}\label{example:1-4}
Let $f: \mathbb{R} \rightarrow \mathbb{R}$ be the function $f(x)=\sin (x)$ and $g: \mathbb{R} \rightarrow \mathbb{R}$ be the function $g(x)=\cos (x)$. Which is even and which is odd?

Even: $\answer{\cos(x)}$ \qquad Odd: $\answer{\sin(x)}$
\end{example}

\begin{exercise}\label{exercise:1-5}
Give an example of a function that is neither even nor odd:
\begin{freeResponse}\end{freeResponse}
\end{exercise}

\begin{exercise}\label{exercise:1-6}
Give an example of a function that is both even and odd:
\begin{freeResponse}\end{freeResponse}
\end{exercise}

Recall that $[a, b]$ is a closed interval of $\mathbb{R}$ and $(a, b)$ is an open interval of $\mathbb{R}$. A function $f$ is called \wordChoice{\choice[correct]{increasing}\choice{decreasing}\choice{constant}\choice{even}} on an interval $I$ if:
\[
f(x)<f(y) \text { for every } x<y \text { in } I .
\]
A function $f$ is called \wordChoice{\choice[correct]{decreasing}\choice{increasing}\choice{constant}\choice{odd}} on an interval $I$ if:
\[
f(x)>f(y) \text { for every } x<y \text { in } I .
\]

Let $f$ and $g$ be two functions. Then
\[
\begin{aligned}
(f+g)(x) & = \answer{f(x)+g(x)} \qquad (f-g)(x) = \answer{f(x)-g(x)} \\
(f g)(x) & = \answer{f(x)g(x)} \\
\left(\frac{f}{g}\right)(x) & = \answer{\frac{f(x)}{g(x)}}
\end{aligned}
\]

\begin{example}\label{example:1-7}
% AUTHOR CHOICE: example chosen during conversion, change freely
Let $f(x)=x^{2}$ and $g(x)=x+1$.
% work: fg = x^2(x+1) = x^3+x^2
\[
\begin{aligned}
(f+g)(x) & = \answer{x^{2}+x+1} \qquad (f-g)(x) = \answer{x^{2}-x-1} \\
(f g)(x) & = \answer{x^{3}+x^{2}} \\
\frac{f}{g}(x) & = \answer{\frac{x^{2}}{x+1}}
\end{aligned}
\]
\end{example}

The \wordChoice{\choice[correct]{composition}\choice{product}\choice{sum}\choice{inverse}} of two functions is the composite function $f \circ g$ where
\[
f \circ g(x)= \answer{f(g(x))}
\]

\begin{warning}
The order matters for composition, just like subtraction and division!
\end{warning}

\begin{example}\label{example:1-8}
% AUTHOR CHOICE: example chosen during conversion, change freely
Let $f(x)=x^{2}$ and $g(x)=x-3$.
\[
(f \circ g)(x) = \answer{(x-3)^{2}} \qquad (g \circ f)(x) = \answer{x^{2}-3}.
\]
\end{example}

\begin{exercise}\label{exercise:1-9}
Try this one:
% AUTHOR CHOICE: example chosen during conversion, change freely
Let $f(x)=2x+1$ and $g(x)=x^{2}$.
\[
(f \circ g)(x)= \answer{2x^{2}+1}
\]
\[
(g \circ f)(x)= \answer{(2x+1)^{2}}.
\]
\end{exercise}
\end{document}
